\section{Introduction}


Cross-chain communication is a fundamental requirement for multi-chain blockchain architectures. As the number of application-specific blockchains (appchains) grows, the ability to transfer assets, share state, and coordinate actions across chains becomes critical for composability and user experience \cite{wood2016polkadot, kwon2016cosmos}.

Existing cross-chain messaging solutions fall into three categories: (i)~trusted bridge committees, which introduce centralization risks \cite{zamyatin2021bridge}; (ii)~relay chains that verify source chain consensus on the destination chain, which are expensive in gas and latency \cite{kwon2016cosmos}; and (iii)~light client protocols that verify block headers, which require maintaining header chains on every connected chain \cite{kiayias2016proofs}.

Lux Warp Messaging (LWM) takes a fourth approach: leveraging the shared validator infrastructure of the Lux Network to produce BLS aggregate signatures that attest to cross-chain messages. Because appchains share validators with the primary network, the destination appchain can verify the source appchain's validator set without maintaining a separate light client. The BLS12-381 pairing-based signature scheme \cite{boneh2001bls, boneh2018compact} enables aggregation of arbitrarily many signatures into a single 48-byte proof, achieving constant-size verification regardless of validator set size.

\subsection{Contributions}

\begin{enumerate}[label=(\roman*)]
    \item We formalize the LWM protocol including message format, signing, aggregation, relay, and verification phases.
    \item We prove security under $2/3$ honest-stake with formal reduction to the co-CDH assumption on BLS12-381.
    \item We introduce optimistic delivery with fraud proofs for latency-sensitive applications.
    \item We provide concrete benchmarks showing sub-second verification and sub-500-byte message overhead.
    \item We analyze relay incentive design to ensure liveness without trusted infrastructure.
\end{enumerate}

\subsection{Paper Organization}

Section~\ref{sec:prelim} covers BLS signature preliminaries and the Lux appchain model. Section~\ref{sec:protocol} specifies the full LWM protocol. Section~\ref{sec:security} provides formal security analysis. Section~\ref{sec:optimistic} introduces the optimistic delivery mode. Section~\ref{sec:relay} discusses relay network design and incentives. Section~\ref{sec:tracking} addresses validator set tracking across appchains. Section~\ref{sec:gas} analyzes on-chain verification costs. Section~\ref{sec:batching} presents message batching for throughput scaling. Section~\ref{sec:comparison} compares LWM with existing protocols. Section~\ref{sec:benchmarks} reports empirical benchmarks. Section~\ref{sec:conclusion} concludes.

